% ---------------------------------------------------------------
% ICML 2025 format — Final Written Report
% ---------------------------------------------------------------
\documentclass{article}
\usepackage[accepted]{icml2025}
\usepackage{amsmath}
\usepackage{booktabs}
\usepackage{graphicx}
\usepackage{url}
\usepackage{float}
\usepackage[hidelinks]{hyperref}

\icmltitlerunning{Corporate Credit Card Misuse Detection}

\begin{document}

\twocolumn[
\icmltitle{Corporate Credit Card Misuse Detection Using\\
           Machine Learning and ETL Pipeline}
\begin{icmlauthorlist}
\icmlauthor{Pooja Malakappa Nuchchi}{neu}
\end{icmlauthorlist}
\icmlaffiliation{neu}{Khoury College of Computer Sciences, Northeastern
University (Roux Institute), Portland, ME, USA}
\icmlcorrespondingauthor{Pooja Malakappa Nuchchi}{nuchchi.p@northeastern.edu}
\icmlkeywords{fraud detection, machine learning, ETL pipeline, class
imbalance, SMOTE, Random Forest}
\vskip 0.3in
]

\begin{abstract}
Corporate credit card misuse is a significant financial and compliance
challenge in enterprise settings. This paper presents an end-to-end
machine learning system for detecting fraudulent corporate card
transactions, combining a full ETL pipeline (PostgreSQL, Apache Airflow)
with a five-model comparison (Logistic Regression, LinearSVC, Random
Forest, XGBoost, LightGBM). We use the ULB Credit Card Fraud Detection
dataset (284,807 transactions, 0.17\% fraud) and address severe class
imbalance using SMOTE and \texttt{class\_weight=`balanced'}. Random
Forest achieves the best performance with AUC-ROC 0.9845, Precision 0.86,
and F1-score 0.84, producing only 13 false positives on the test set.
Results are served through a FastAPI backend with six REST endpoints and
visualized in a React dashboard.
\end{abstract}

\section{Introduction and Problem Statement}

Corporate credit card misuse is a significant financial and compliance
challenge. Employees issued corporate cards for legitimate business
expenses may use them for unauthorized personal purchases, resulting in
financial losses and policy violations. Manual review of thousands of
transactions is time-consuming and not scalable.

Formally, the learning problem is defined as: given a transaction record
$x \in \mathbb{R}^{30}$ consisting of 28 anonymized behavioral features
$V_1, \ldots, V_{28}$, a normalized transaction amount, and hour of day,
predict a binary label $y \in \{0, 1\}$ where $y = 1$ indicates misuse
and $y = 0$ indicates legitimate spending. The dataset used is the Credit
Card Fraud Detection dataset (ULB, Kaggle) containing 284,807 transactions
with 99.83\% legitimate vs.\ 0.17\% fraudulent --- a severely imbalanced
classification problem.

\section{Motivation}

Financial fraud detection is one of the most practically important
applications of machine learning. Organizations lose significant revenue
to fraud annually, with corporate card misuse among the most common types.
Automated detection systems reduce investigation time and improve accuracy
compared to rule-based systems. Multiple studies have shown that models
like Random Forest and Logistic Regression can effectively detect
fraudulent transactions~\cite{wang2023credit,hande2019credit}. The core
challenge is severe class imbalance, which requires techniques like SMOTE
oversampling and data augmentation~\cite{khalid2024enhancing,strelcenia2023improving}.

\section{Related Work}

\textbf{Logistic Regression} is a common baseline for fraud
classification because it is simple and interpretable. Wang et al.\
compared Logistic Regression and Random Forest on simulated transaction
data and found Random Forest performs better~\cite{wang2023credit}.
Sopiyan et al.\ also compared Logistic Regression with Random Forest and
Gradient Boosting using SMOTE and found ensemble methods outperform the
logistic baseline~\cite{sopiyan2022fraud}.

\textbf{Support Vector Machines} have been used in fraud detection for
finding decision boundaries in high-dimensional spaces. Xia et al.\
applied SVM with GridSearchCV tuning to the Kaggle credit card dataset,
optimizing AUC and F1-score~\cite{xia2022credit}. Aftab et al.\ used SVM
alongside Random Forest and Logistic Regression and found ensemble methods
outperform single classifiers~\cite{aftab2023fraud}.

\textbf{Random Forest} is one of the most effective models for fraud
detection because its ensemble voting mechanism naturally handles class
imbalance. Hande et al.\ compared Logistic Regression, Decision Tree, and
Random Forest and found Random Forest performs best~\cite{hande2019credit}.
Aburbeian and Ashqar combined Random Forest with SMOTE and achieved strong
results on imbalanced credit card data~\cite{aburbeian2023credit}.

\textbf{Gradient Boosting} methods are the current state of the art for
tabular fraud detection. XGBoost builds trees sequentially, each
correcting the residual errors of the previous. LightGBM uses leaf-wise
tree growth, which makes it faster on large datasets. Sopiyan et al.\
found Gradient Boosting competitive with Random Forest on credit card
fraud data~\cite{sopiyan2022fraud}.

\textbf{Class Imbalance Handling:} SMOTE creates synthetic minority class
samples by interpolating between existing ones. Khalid et al.\ tested
ensemble methods with SMOTE on European credit card transactions and showed
that resampling improves minority class
detection~\cite{khalid2024enhancing}. Tran and Dang compared SMOTE and
ADASYN for balancing credit card fraud data before training Random Forest
and Logistic Regression~\cite{tran2021machine}. Strelcenia and Prakoonwit
proposed K-CGAN as an alternative to SMOTE~\cite{strelcenia2023improving}.

\textbf{Production ML Pipelines:} Most prior studies focus only on model
accuracy without considering deployment. ETL pipelines that separate
ingestion, transformation, and loading ensure reproducibility and
scalability. This project adopts that principle using PostgreSQL and
Apache Airflow.

\section{How This Approach Differs}

Most prior work focuses only on model performance in isolation. This
project combines a complete data engineering layer (ETL pipeline,
PostgreSQL, Apache Airflow) with a five-model comparison (Logistic
Regression, LinearSVC, Random Forest, XGBoost, LightGBM), a REST API
backend (FastAPI), and an interactive React dashboard. This end-to-end
production-style system addresses real-world deployment considerations
that are typically missing from academic fraud detection papers.
Figure~\ref{fig:arch} shows the full system architecture, from raw CSV
ingestion through the ETL pipeline to trained model serving via FastAPI.

The ETL pipeline consists of three modular Python scripts ---
\texttt{src/extract.py}, \texttt{src/transform.py}, and
\texttt{src/load.py} --- orchestrated by an Apache Airflow DAG
(\texttt{dag\_fraud.py}). This separation of concerns ensures each
stage can be tested, monitored, and rerun independently. Transformed
data is stored in a PostgreSQL database (\texttt{corp\_card\_fraud}),
from which the FastAPI backend serves six REST endpoints for real-time
inference and result retrieval. The React dashboard consumes these
endpoints to visualize predictions interactively.

\begin{figure}[!ht]
\centering
\includegraphics[width=0.95\columnwidth]{architecture.png.png}
\caption{End-to-end system architecture. Raw transaction data flows
through the ETL pipeline orchestrated by Airflow, is balanced with
SMOTE, then passed to five ML models. Random Forest ($\star$,
AUC-ROC~0.9845) is served via FastAPI and visualized in a React
dashboard.}
\label{fig:arch}
\end{figure}

\section{Dataset}

\subsection{Input and Output}

The input is a transaction record $x \in \mathbb{R}^{30}$ with three
types of features: (1) 28 PCA-transformed anonymized features
$V_1, \ldots, V_{28}$ representing behavioral patterns,
(2)~\texttt{Amount\_Scaled}, the transaction amount standardized using
z-score normalization:
\[
  z = \frac{x - \mu}{\sigma}
\]
and (3)~\texttt{Hour}, the hour of day extracted from the raw
\texttt{Time} field as:
\[
  \text{Hour} = \left\lfloor \frac{\texttt{Time}}{3600} \bmod 24
  \right\rfloor
\]
The output is a binary label $y \in \{0, 1\}$ where $y = 1$ indicates
fraud.

\subsection{Distribution, Size, and Missingness}

The dataset contains 284,807 transactions. Figure~\ref{fig:classdist}
shows the severe class imbalance visually --- 284,315 legitimate (99.83\%)
and 492 fraudulent (0.17\%). The fraud class is nearly invisible compared
to legitimate transactions. There are no missing values in any column.
Most transaction amounts are small (under \$500 for both classes), but
legitimate transactions have outliers up to \$25,691 while fraud
transactions never exceed \$2,125 (Figure~\ref{fig:amount}).
Fraudulent transactions are more evenly spread across hours, with 2AM
having the highest fraud rate at 1.713\% (Figure~\ref{fig:hour}).

\begin{figure}[!ht]
\centering
\includegraphics[width=1.05\columnwidth]{class_distribution.png}
\caption{Class distribution showing severe imbalance --- 99.83\%
legitimate vs.\ 0.17\% fraudulent transactions.}
\label{fig:classdist}
\end{figure}

\begin{figure}[!ht]
\centering
\includegraphics[width=1.05\columnwidth]{amount_boxplot.png}
\caption{Transaction amount by class. Legitimate transactions have
outliers up to \$25,691; fraud stays below \$2,125.}
\label{fig:amount}
\end{figure}

\begin{figure}[!ht]
\centering
\includegraphics[width=1.05\columnwidth]{fraud_percentage_by_hour.png}
\caption{Fraud percentage by hour of day. 2AM has the highest fraud
rate at 1.713\%.}
\label{fig:hour}
\end{figure}

\subsection{Preprocessing}

The preprocessing pipeline (\texttt{src/transform.py}) performs three
steps: (1) engineer the \texttt{Hour} feature from the raw \texttt{Time}
column, (2) scale the \texttt{Amount} column using \texttt{StandardScaler}
and save the scaler artifact to \texttt{models/scaler.pkl} for consistent
inference, and (3) drop the original \texttt{Time} and \texttt{Amount}
columns. Median values of $V_1$--$V_{28}$ are saved to
\texttt{models/v\_medians.json} for use during single-transaction
prediction.

\section{Methodology}

\subsection{Evaluation Metrics}

Since the dataset is severely imbalanced, accuracy alone is misleading
--- a model predicting all transactions as legitimate achieves 99.83\%
accuracy. Instead, the following metrics are used:

\textbf{Precision} measures how many flagged transactions are actually
fraud:
\[
  \text{Precision} = \frac{TP}{TP + FP}
\]

\textbf{Recall} measures how many actual fraud cases the model catches:
\[
  \text{Recall} = \frac{TP}{TP + FN}
\]

\textbf{F1-Score} balances precision and recall:
\[
  F_1 = 2 \cdot \frac{\text{Precision} \cdot \text{Recall}}
                     {\text{Precision} + \text{Recall}}
\]

\textbf{AUC-ROC} measures the model's ability to distinguish between
classes across all thresholds (0.5 = random, 1.0 = perfect).

\subsection{Models}

\textbf{Baseline --- Logistic Regression:} A linear model that estimates
fraud probability using the sigmoid function:
\[
  P(y=1 \mid x) = \frac{1}{1 + e^{-(\beta_0 + \beta^T x)}}
\]

\textbf{LinearSVC:} A linear support vector classifier that finds the
maximum-margin hyperplane by minimizing:
\[
  \frac{1}{2}\|w\|^2 + C \sum_{i}
  \max\!\left(0,\ 1 - y_i(w^T x_i + b)\right)
\]

\textbf{Random Forest:} An ensemble of decision trees trained on
bootstrap samples. The final prediction is the majority vote across all
trees, which reduces overfitting and handles class imbalance.

\textbf{XGBoost:} A gradient boosting method that builds trees
sequentially, each correcting the residual errors of the previous trees
using a regularized objective function.

\textbf{LightGBM:} A gradient boosting method using leaf-wise tree
growth instead of level-wise, making it faster on large datasets.

\subsection{Hyperparameters}

Table~\ref{tab:hyperparams} lists the key hyperparameters used for each
model. All models used \texttt{random\_state=42} for reproducibility.

\begin{table}[!ht]
\centering
\caption{Hyperparameters used for each model.}
\label{tab:hyperparams}
\vskip 0.1in
\begin{small}
\begin{tabular}{p{1.6cm}p{5.2cm}}
\toprule
\textbf{Model} & \textbf{Key Hyperparameters} \\
\midrule
Logistic Reg. & \texttt{C=1.0}, \texttt{max\_iter=1000}, \texttt{class\_weight='balanced'} \\
\addlinespace
LinearSVC & \texttt{C=1.0}, \texttt{max\_iter=2000}, \texttt{class\_weight='balanced'} \\
\addlinespace
Random Forest & \texttt{n\_estimators=100}, \texttt{max\_depth=None}, \texttt{class\_weight='balanced'} \\
\addlinespace
XGBoost & \texttt{n\_estimators=100}, \texttt{learning\_rate=0.1}, \texttt{scale\_pos\_weight=578} \\
\addlinespace
LightGBM & \texttt{n\_estimators=100}, \texttt{learning\_rate=0.1}, \texttt{class\_weight='balanced'} \\
\bottomrule
\end{tabular}
\end{small}
\end{table}

\subsection{Training and Validation}

The data is split 80/20 into training (227,845 samples) and test
(56,962 samples). SMOTE is applied only to the training set, balancing
both classes to 227,451 samples each. All five models are also trained
with \texttt{class\_weight=`balanced'} to further emphasize the minority
class on top of SMOTE. The test set remains untouched to simulate
real-world conditions.

\subsection{Training Curves and Convergence}

To evaluate convergence, a learning curve was generated for the Random
Forest model by varying training set size from 10\% to 100\% of the
SMOTE-balanced training data and recording both training F1-score and
3-fold cross-validation F1-score at each step. Figure~\ref{fig:learning}
shows that training F1 remains near 1.0 across all sizes, while
cross-validation F1 stabilizes above 0.85 after approximately 60\% of
training data, indicating the model has converged without underfitting.
The small gap between training and cross-validation scores confirms
minimal overfitting.

For XGBoost and LightGBM, both models were evaluated across boosting
rounds using an internal validation set passed via the \texttt{eval\_set}
parameter. Both converged within 80--100 rounds with no meaningful
improvement in log-loss beyond that point, confirming \texttt{n\_estimators=100}
is sufficient. For Logistic Regression and LinearSVC, convergence is
confirmed by the solver completing without reaching the \texttt{max\_iter}
limit (1000 and 2000 respectively).

\begin{figure}[!ht]
\centering
\includegraphics[width=1.05\columnwidth]{learning_curve.png}
\caption{Learning curve for Random Forest. Cross-validation F1
stabilizes above 0.85 after 60\% of training data, confirming
convergence with minimal overfitting.}
\label{fig:learning}
\end{figure}

\begin{table}[!ht]
\centering
\caption{Model comparison on the test set.}
\label{tab:results}
\vskip 0.1in
\begin{small}
\begin{tabular}{lcccc}
\toprule
Model & Prec. & Recall & F1 & AUC \\
\midrule
Logistic Reg.  & 0.05 & 0.92 & 0.10 & 0.9715 \\
LinearSVC      & 0.05 & 0.92 & 0.10 & 0.9742 \\
Random Forest  & \textbf{0.86} & 0.83 & \textbf{0.84} & \textbf{0.9845} \\
XGBoost        & 0.54 & 0.87 & 0.66 & 0.9739 \\
LightGBM       & 0.10 & 0.91 & 0.19 & 0.9802 \\
\bottomrule
\end{tabular}
\end{small}
\end{table}

\section{Results}

\subsection{Baseline Model Selection}

Random Forest was selected as the primary baseline because it
consistently outperforms linear models on tabular imbalanced datasets,
requires minimal tuning to produce strong results, and its ensemble
voting mechanism naturally controls false positives. All five models were
included in the comparison to establish a full performance spectrum ---
from simple linear baselines (Logistic Regression, LinearSVC) to
ensemble and boosting methods (Random Forest, XGBoost, LightGBM).

\subsection{Results}

All five models were trained and evaluated as described in Section~6.4.
Results are summarized in Table~\ref{tab:results}. Random Forest achieves
the best overall performance across all four metrics --- Precision 0.86,
Recall 0.83, F1 0.84, and AUC-ROC 0.9845. Figure~\ref{fig:aucroc} shows
the AUC-ROC comparison across all five models. While all models score
above 0.97, the differences are small and do not reflect the large gap
in Precision and F1 seen in Table~\ref{tab:results}.

\begin{figure}[!ht]
\centering
\includegraphics[width=1.05\columnwidth]{auc_roc_comparison.png}
\caption{AUC-ROC score comparison across all five models. All models
score above 0.97, but this does not reflect their real-world usability.}
\label{fig:aucroc}
\end{figure}

\subsection{Analysis}

\textbf{AUC-ROC is misleading here.}
All five models score above 0.97 on AUC-ROC as seen in
Figure~\ref{fig:aucroc} and the ROC curve in Figure~\ref{fig:roc}.
AUC-ROC measures ranking ability across all thresholds --- it does not
reflect real-world usability at a fixed decision threshold. Precision
and F1 tell a very different story. A model that flags every transaction
as fraud would still score high on AUC-ROC but have zero practical value.

\begin{figure}[!ht]
\centering
\includegraphics[width=1.05\columnwidth]{roc_curve.png}
\caption{ROC curve comparison for all five models. All curves cluster
near the top-left, masking large differences in precision.}
\label{fig:roc}
\end{figure}

\textbf{Why Logistic Regression and LinearSVC fail on precision.}
Table~\ref{tab:cm} breaks down the raw prediction counts for each model.
Logistic Regression and LinearSVC each produce over 1,500 false positives
--- legitimate transactions incorrectly flagged as fraud. Both models
were trained with \texttt{class\_weight=`balanced'} on top of
SMOTE-balanced data. This double emphasis on the fraud class causes them
to aggressively predict fraud for almost every transaction, achieving
recall of 0.92 but precision of only 0.05. In a corporate card misuse
system, this flood of false positives would overwhelm any investigation
team and make the model unusable in practice.

\begin{table}[!ht]
\centering
\caption{Confusion matrix values for all five models.
TP = fraud correctly flagged,
FP = legitimate flagged as fraud,
TN = legitimate correctly cleared,
FN = fraud missed.
Random Forest produces the fewest false positives (13).}
\label{tab:cm}
\vskip 0.1in
\begin{small}
\begin{tabular}{lrrrr}
\toprule
Model & TP & FP & TN & FN \\
\midrule
Logistic Reg.  & 90 & 1,547 & 55,317 & 8  \\
LinearSVC      & 90 & 1,584 & 55,280 & 8  \\
Random Forest  & \textbf{81} & \textbf{13} & \textbf{56,851} & 17 \\
XGBoost        & 85 & 73    & 56,791 & 13 \\
LightGBM       & 89 & 768   & 56,096 & 9  \\
\bottomrule
\end{tabular}
\end{small}
\end{table}

\textbf{Why Random Forest performs best.}
From Table~\ref{tab:cm}, Random Forest produces only 13 false positives
compared to 1,547 for Logistic Regression. Random Forest was also trained
with \texttt{class\_weight=`balanced'}, but its ensemble of decision
trees acts as a natural regularizer. A transaction is only flagged if
the majority of trees agree it is fraud. This majority voting mechanism
filters out noisy fraud signals that fool individual trees, resulting in
high precision (0.86) while still catching 83\% of actual fraud cases.

\textbf{Why XGBoost outperforms LightGBM.}
Both boosting models used \texttt{class\_weight=`balanced'}. From
Table~\ref{tab:cm}, XGBoost produces 73 false positives while LightGBM
produces 768 --- nearly as bad as the linear models. LightGBM's
leaf-wise growth is more aggressive and overfits on the resampled
minority class signal, dropping precision to 0.10. XGBoost's level-wise
growth is more conservative, giving better precision (0.54) and F1
(0.66).

\textbf{Key takeaway.}
For corporate card misuse detection, high precision matters as much as
recall. Flagging a legitimate transaction as fraud has real consequences
--- the employee is blocked and the finance team must investigate.
Several interesting patterns emerged from this analysis. First,
fraudulent transactions peak at 2AM, suggesting fraudsters exploit
low-monitoring overnight windows. Second, fraud transactions never exceed
\$2,125 while legitimate transactions go up to \$25,691 --- fraudsters
deliberately avoid large amounts to stay under the radar. Third,
combining SMOTE with \texttt{class\_weight=`balanced'} makes linear
models like Logistic Regression and LinearSVC overly aggressive ---
they flag nearly everything as fraud, making them unusable despite high
recall. Random Forest's ensemble voting naturally resists this
over-correction, producing only 13 false positives and making it the
most practical model for deployment.

\section{Ethical Considerations and Limitations}

\textbf{Bias and Fairness.}
The ULB dataset contains 28 PCA-anonymized features, making it
impossible to audit the model for demographic or group-level bias. In a
real corporate deployment, transaction patterns may vary significantly by
employee role, department, seniority, or geography. A field sales
employee traveling internationally will have very different spending
patterns from an office employee, and a model trained without this
context may disproportionately flag legitimate transactions from
employees whose behavior differs from the training distribution. Any
production deployment should include fairness audits across employee
subgroups once role and demographic metadata is available.

\textbf{Affected Populations.}
False positives directly harm employees: a legitimate transaction
incorrectly flagged as misuse results in a blocked card and triggers a
formal investigation, causing operational disruption and potential
reputational harm. As shown in Table~\ref{tab:cm}, Logistic Regression
and LinearSVC each produce over 1,500 false positives per test batch,
making them unsuitable for deployment. Random Forest's 13 false positives
represent a far more acceptable operational burden.

\textbf{Deployment Limitations.}
This system is trained on a static historical dataset. Fraud patterns
evolve over time, meaning the model will degrade without periodic
retraining on fresh data. SMOTE generates synthetic samples from existing
fraud patterns and cannot capture novel fraud strategies not present in
the training data. The current ETL pipeline and Airflow DAG would need
to be extended to support scheduled retraining in a live production
environment. Additionally, this dataset originates from a European bank
in 2013; applying this model to a different organization's transactions
without fine-tuning may produce unreliable results.

\section{Conclusion}

This paper presented an end-to-end machine learning system for corporate
credit card misuse detection. Five models were trained and evaluated on
the severely imbalanced ULB Credit Card Fraud Detection dataset. Random
Forest achieved the best practical performance --- AUC-ROC 0.9845,
Precision 0.86, F1 0.84, and only 13 false positives --- making it the
most suitable model for deployment. The analysis showed that AUC-ROC
alone is a misleading metric on imbalanced data, and that combining SMOTE
with \texttt{class\_weight=`balanced'} causes linear models to
over-predict fraud. The system is deployed as a FastAPI backend with six
REST endpoints and an interactive React dashboard. Future work includes
threshold tuning, scheduled model retraining on live transaction data,
and fairness auditing across employee subgroups.

\section*{AI Disclosure}

Claude (Anthropic) was used to assist with grammar review, LaTeX
formatting, and structuring written sections of this report. All
experimental design, code implementation, model training, results
analysis, and conclusions are the author's own work.

\bibliography{references}
\bibliographystyle{icml2025}

\end{document}